% ============================================================================
% Fichier  : partie0/P00_decision_mairie.tex
% Titre    : Prologue : la Décision de la Mairie
% Auteur   : MINKA MI NGUIDJOÏ Thierry Emmanuel
% ============================================================================
% Fil rouge : Mairie d'Eséka (Nyong-et-Kéllé). Cas fictif ; mécanismes et
% références normatives exacts.
% ============================================================================

\chapter{Prologue : la Décision de la Mairie}
\label{chap:prologue}
\chaptermark{Prologue}

\epigraph{%
    « La sécurité n'est pas un produit que l'on achète,\\
    c'est une décision que l'on tient dans la durée. »%
}{Principe du cours}

\niveauChapitre{Introductif}{Aucun prérequis. Les notions sont reprises en
détail dans les chapitres.}

\begin{boxobjectifs}
\begin{itemize}
    \item Comprendre qu'un incident de sécurité est d'abord un problème de
          management, non de seule technique.
    \item Distinguer la réaction ponctuelle d'une démarche durable et pilotée.
    \item Découvrir le fil rouge de la mairie d'Eséka et ses quatre phases.
    \item Situer le cycle \sigle{PDCA} au cœur du cours.
    \item Formuler les questions auxquelles le cours répond.
\end{itemize}
\end{boxobjectifs}

\bigskip

% ============================================================================
\section*{13 janvier 2026, 7 h 50, mairie d'Eséka}
% ============================================================================

L'agent d'état civil arrive avant l'ouverture au public. Il branche, comme
chaque matin, la clé \texttt{USB} sur laquelle il a travaillé la veille chez
lui. En quelques minutes, les fichiers du poste deviennent illisibles~: les
actes de naissance numérisés, le registre des mariages, le tableur des
redevances ne s'ouvrent plus. À 8~h~15, trois postes du service sont touchés,
reliés au même dossier partagé. À 9~h, une file d'usagers attend des copies
d'actes que la mairie ne peut plus délivrer.

\medskip
Ce que personne ne sait encore ce matin-là\footnote{Cas fictif établi pour le
cours. Les noms (mairie d'Eséka, cabinet \textsc{Matomb Conseil}, agents
cités), les dates et les montants sont fictifs~; les mécanismes de gestion et
les références normatives (ISO/IEC 27001, 27002, 27005, lois camerounaises)
sont exacts.}~: le mal n'est pas la clé \texttt{USB}. Le mal, c'est qu'un seul
poste suffise à paralyser un service, qu'aucune sauvegarde récente n'existe,
que le même mot de passe serve à tout le monde, et que personne, dans la
mairie, n'ait jamais décidé qui est responsable de quoi en matière de
sécurité.

% ============================================================================
\section{Réparer ne suffit pas}
\label{sec:p00-reparer}
% ============================================================================

Le \sigle{DSI} de la commune, M.~\textsc{Biloa}, rétablit les postes en deux
jours à partir d'une vieille copie. Le service rouvre. L'affaire pourrait
s'arrêter là~: on a réparé. Mais le maire, M.~\textsc{Essama}, pose la bonne
question en conseil~: « Qu'est-ce qui nous garantit que cela ne
recommencera pas la semaine prochaine, en pire~? »

\medskip
La réponse est désagréable~: rien. Réparer traite l'incident du jour~; cela ne
traite pas la \termecle{cause}. La cause est organisationnelle~: pas de
politique, pas de rôles, pas d'appréciation des risques, pas de règles de
sauvegarde, pas de contrôle. Tant que ces manques demeurent, le prochain
incident est une question de temps. La mairie n'a pas un problème technique~:
elle a un problème de \termecle{management de la sécurité}.

\begin{boxmetier}{RSSI}
\textbf{Du pompier au pilote.} Réagir à chaque incident, c'est être pompier.
Un \sigle{RSSI} construit un système qui réduit le nombre d'incendies et
limite leurs dégâts. Il ne promet pas le risque zéro~: il promet un risque
\emph{connu}, \emph{décidé} et \emph{maîtrisé}.
\end{boxmetier}

% ============================================================================
\section{La décision : un système, pas un achat}
\label{sec:p00-decision}
% ============================================================================

Le maire ne débloque pas un budget pour « acheter de la sécurité ». Il prend
une décision de gouvernance~: doter la mairie d'un \termecle{système de
management de la sécurité de l'information} (\sigle{SMSI}), c'est-à-dire un
ensemble de règles, de rôles et d'activités qui, ensemble, maintiennent la
sécurité dans la durée. Il nomme Mme~\textsc{Abena} \sigle{RSSI} de la commune
et mandate le cabinet \textsc{Matomb Conseil} pour l'accompagner vers une
certification \norme{ISO/IEC 27001}.

\begin{boxdef}
\termecle{Système de management de la sécurité de l'information} (\sigle{SMSI})~:
partie du système de management global qui, sur la base d'une appréciation des
risques, établit, met en \oe uvre, exploite, surveille et améliore la sécurité
de l'information~\cite{iso27001}. Ce n'est pas un logiciel~: c'est une manière
organisée de décider et de tenir ses décisions.
\end{boxdef}

% ============================================================================
\section{Le fil rouge : quatre phases}
\label{sec:p00-filrouge}
% ============================================================================

Le cours suit la mairie d'Eséka, dans le département du Nyong-et-Kéllé, tout au
long de sa démarche. Quatre phases, nommées d'après la géographie locale,
structurent le fil rouge et reviennent dans les chapitres.

\begin{tableaucours}{L{3.2cm}L{9.0cm}}{%
    \tableauHeader{Phase} &
    \tableauHeader{Ce qu'elle installe}
}
Opération \textsc{Nyong} & La gouvernance et la politique de sécurité (parties I et II) \\
Opération \textsc{Kellé} & L'appréciation et le traitement des risques (partie IV) \\
Opération \textsc{Makak} & La mise en \oe uvre des mesures (partie V) \\
Opération \textsc{Matomb} & L'audit interne et la certification (partie VI) \\
\end{tableaucours}
\vspace{-0.8em}
\begin{center}{\small\itshape Tableau~0.1. Les quatre phases du fil rouge}\end{center}

% ============================================================================
\section{Le cœur du cours : le cycle PDCA}
\label{sec:p00-pdca}
% ============================================================================

Un \sigle{SMSI} n'est pas un projet qui se termine~: c'est une \termecle{roue}
qui tourne. Le cycle \sigle{PDCA} (\emph{Plan, Do, Check, Act}), ou roue de
Deming, en est le moteur et revient à chaque partie du cours.

\begin{tableaucours}{L{2.6cm}L{9.6cm}}{%
    \tableauHeader{Étape} &
    \tableauHeader{Ce que fait la mairie}
}
Planifier (\emph{Plan}) & Apprécier les risques, fixer la politique et les objectifs \\
Déployer (\emph{Do}) & Mettre en \oe uvre les mesures et former les agents \\
Contrôler (\emph{Check}) & Mesurer, auditer, passer en revue de direction \\
Agir (\emph{Act}) & Corriger, améliorer, décider le cycle suivant \\
\end{tableaucours}
\vspace{-0.8em}
\begin{center}{\small\itshape Tableau~0.2. Le cycle PDCA appliqué à la mairie}\end{center}

\begin{boxinfo}
\textbf{Les questions auxquelles ce cours répond.}
\begin{itemize}
    \item Qui décide de la sécurité, et au nom de quoi~?
    \item Comment connaître ses risques avant qu'ils ne deviennent des
          incidents~?
    \item Quelles mesures mettre en \oe uvre, et comment prouver qu'elles
          fonctionnent~?
    \item Comment faire reconnaître, par un audit, que la sécurité est
          maîtrisée~?
\end{itemize}
\end{boxinfo}

\separateur

\begin{boxresume}
\textbf{Ce qu'il faut retenir de ce prologue~:}
\begin{itemize}
    \item Un incident est d'abord un révélateur de \textbf{causes
          organisationnelles}, pas seulement un problème technique.
    \item \textbf{Réparer} traite le symptôme~; \textbf{manager} traite la
          cause.
    \item Un \textbf{SMSI} est un système durable de décisions tenues, non un
          achat ponctuel.
    \item Le fil rouge suit la \textbf{mairie d'Eséka} en quatre phases.
    \item Le \textbf{cycle PDCA} est le moteur du cours.
\end{itemize}
\end{boxresume}

\begin{boxinfo}
\textbf{Transition.} Avant de manager la sécurité, il faut la définir~: le
Chapitre~\ref{chap:cia} pose ce qu'est la sécurité de l'information.
\end{boxinfo}
